% ==============================================================================
%  5.2  Estimación de los coeficientes de regresión. La matriz hat
% ==============================================================================
\section{Estimación de los coeficientes de regresión}
\label{sec:t5-estimacion}

El modelo se ajusta por el método de mínimos cuadrados (\cref{sec:t2-mc}). Para un valor dado de
$\vect{\beta}$, los errores son $\vect{\varepsilon}=\vect{Y}-\vect{X}\vect{\beta}$, y la suma de
sus cuadrados se escribe como un producto escalar:
\[
  \sumi \varepsilon_i^2 =
  \begin{pmatrix} \varepsilon_1 & \varepsilon_2 & \cdots & \varepsilon_n \end{pmatrix}
  \begin{pmatrix} \varepsilon_1\\ \varepsilon_2\\ \vdots\\ \varepsilon_n \end{pmatrix}
  = \vect{\varepsilon}\T\vect{\varepsilon} .
\]
Se trata, por tanto, de minimizar en $\vect{\beta}$
\[
  \vect{\varepsilon}\T\vect{\varepsilon} = (\vect{Y}-\vect{X}\vect{\beta})\T(\vect{Y}-\vect{X}\vect{\beta}) .
\]
Derivando respecto de $\vect{\beta}$,
\[
  \frac{\partial}{\partial\vect{\beta}}
  \Big((\vect{Y}-\vect{X}\vect{\beta})\T(\vect{Y}-\vect{X}\vect{\beta})\Big)
  = -2\vect{X}\T(\vect{Y}-\vect{X}\vect{\beta}),
\]
e igualando a $\vect{0}$ se obtienen las \textbf{ecuaciones normales}.

\begin{teorema}[Estimador de mínimos cuadrados][teo:t5-mc]
  Las ecuaciones normales son
  \[
    (\XtX)\,\bhat = \vect{X}\T\vect{Y},
  \]
  y, si $\XtX$ es invertible, su solución es
  \[
    \boxed{\ \bhat = \XtXi\vect{X}\T\vect{Y}\ } .
  \]
\end{teorema}

\begin{nota}
  $\XtX$ es invertible si las columnas de $\vect{X}$ son linealmente independientes; en la
  regresión simple, si no todos los $X_i$ son iguales ($\SSX>0$).
\end{nota}

\subsection{Valores ajustados, residuos y matriz hat}

Una vez estimados los coeficientes se obtienen el vector de \textbf{valores ajustados} y el
vector de \textbf{residuos} (\cref{def:t2-ajuste}):
\[
  \est{\vect{Y}}_{n\times1} = \begin{pmatrix} \est{Y}_1\\ \vdots\\ \est{Y}_n \end{pmatrix},
  \qquad
  \vect{e}_{n\times1} = \begin{pmatrix} e_1\\ \vdots\\ e_n \end{pmatrix}.
\]
Los valores ajustados se calculan como
\[
  \est{\vect{Y}} = \vect{X}\bhat = \vect{X}\XtXi\vect{X}\T\vect{Y},
\]
y los residuos $e_i=Y_i-\est{Y}_i$ se escriben matricialmente como
\[
  \vect{e} = \vect{Y}-\est{\vect{Y}} = \vect{Y}-\vect{X}\bhat
  = \vect{Y}-\vect{X}\XtXi\vect{X}\T\vect{Y}
  = \big(\vect{I}-\vect{X}\XtXi\vect{X}\T\big)\vect{Y} .
\]

\begin{definicion}[Matriz hat][def:t5-hat]
  La \textbf{matriz hat} es
  \[
    \vect{H}_{n\times n} = \vect{X}\XtXi\vect{X}\T \in\R^{n\times n} .
  \]
  Con ella,
  \[
    \est{\vect{Y}} = \vect{H}\vect{Y}, \qquad \vect{e} = (\vect{I}-\vect{H})\vect{Y} .
  \]
\end{definicion}

La matriz hat desempeña un papel muy importante en la identificación de puntos influyentes: su
elemento $(i,k)$ es el coeficiente $h_{ik}$ de $Y_k$ en el valor ajustado $\est{Y}_i$
(\cref{prop:t4-hik}), y su diagonal contiene los leverages $h_{ii}$ (\cref{def:t4-leverage}).

\begin{proposicion}[Propiedades de la matriz hat][prop:t5-hat]
  La matriz hat es \textbf{simétrica} e \textbf{idempotente}:
  \[
    \vect{H}\T=\vect{H}, \qquad \vect{H}\vect{H}=\vect{H},
  \]
  y por tanto $\vect{H}\T\vect{H}=\vect{H}\vect{H}=\vect{H}$. La matriz de varianzas-covarianzas de
  los residuos es
  \[
    \Var(\vect{e}) = \sigma^2(\vect{I}-\vect{H}) .
  \]
\end{proposicion}

\subsection{Propiedades de los residuos en forma matricial}

El modelo impone que los errores $\varepsilon_i$ sean independientes, normales y con la misma
varianza. Los errores no son observables; los residuos $e_i$ sí lo son, y se utilizan como
estimadores de los $\varepsilon_i$. Sin embargo, los residuos no son independientes entre sí y sus
varianzas no son iguales (\cref{prop:t4-prop-residuos}):
\begin{itemize}
  \item de $\Var(\vect{e})=\sigma^2(\vect{I}-\vect{H})$, la varianza de $e_i$ es el elemento $i$ de
    la diagonal,
    \[
      \Var(e_i) = \sigma^2(1-h_{ii}),
    \]
    que depende del $i$-ésimo elemento de la diagonal de la matriz hat, el leverage de la
    observación $i$, que se utiliza para detectar observaciones atípicas o potencialmente
    influyentes (\cref{sec:t4-outliers-x});
  \item la covarianza entre $e_i$ y $e_j$ es el elemento $(i,j)$ de $\sigma^2(\vect{I}-\vect{H})$,
    \[
      \Cov(e_i,e_j) = \sigma^2(-h_{ij}) = -\sigma^2h_{ij}, \qquad i\neq j .
    \]
\end{itemize}
Se tendría $\Var(e_i)=\sigma^2$ y $\Cov(e_i,e_j)=0$ solo si $\vect{H}=\vect{0}$, lo que no tiene
sentido, ya que entonces $\est{\vect{Y}}=\vect{H}\vect{Y}=\vect{0}$.

\subsection{El caso de la regresión simple}

\begin{proposicion}[Estimadores y matriz hat en la regresión simple][prop:t5-simple]
  Con $\vect{X}$ la matriz de diseño de la \cref{def:t5-modelo},
  \[
    \XtX = \begin{pmatrix} n & \sumi X_i\\ \sumi X_i & \sumi X_i^2 \end{pmatrix},
    \qquad
    \XtXi = \begin{pmatrix}
      \dfrac1n+\dfrac{\media{X}^2}{\SSX} & -\dfrac{\media{X}}{\SSX}\\[10pt]
      -\dfrac{\media{X}}{\SSX} & \dfrac{1}{\SSX}
    \end{pmatrix},
  \]
  y se obtienen:
  \begin{enumerate}
    \item los estimadores de mínimos cuadrados del Tema 2 (\cref{teo:t2-estimadores-mc}),
      $\est{\beta}_1=SS_{XY}/\SSX$ y $\est{\beta}_0=\media{Y}-\est{\beta}_1\media{X}$;
    \item los elementos de la matriz hat del Tema 4 (\cref{prop:t4-hik}),
      $h_{ik}=\frac1n+\frac{(X_i-\media{X})(X_k-\media{X})}{\SSX}$.
  \end{enumerate}
\end{proposicion}

\begin{demostracion}
  \emph{Inversa de $\XtX$.} Se usan las identidades
  \[
    \sumi X_i = n\media{X},
    \qquad
    \sumi X_i^2 = \SSX + n\media{X}^2 .
  \]
  El determinante de $\XtX$ es
  \begin{align*}
    \det(\XtX) &= n\sumi X_i^2-\left(\sumi X_i\right)^2 \\
    &= n\left(\SSX+n\media{X}^2\right) - n^2\media{X}^2 \\
    &= n\SSX ,
  \end{align*}
  así que
  \begin{align*}
    \XtXi &= \frac{1}{n\SSX}\begin{pmatrix} \sumi X_i^2 & -\sumi X_i\\[4pt] -\sumi X_i & n \end{pmatrix} \\[4pt]
    &= \frac{1}{n\SSX}\begin{pmatrix} \SSX+n\media{X}^2 & -n\media{X}\\ -n\media{X} & n \end{pmatrix} ,
  \end{align*}
  que, dividiendo cada elemento por $n\SSX$, es la expresión del enunciado.
  \begin{enumerate}
    \item Como
      \[
        \vect{X}\T\vect{Y} = \begin{pmatrix} \sumi Y_i\\[4pt] \sumi X_iY_i \end{pmatrix}
        = \begin{pmatrix} n\media{Y}\\[4pt] \sumi X_iY_i \end{pmatrix},
      \]
      el producto $\bhat=\XtXi\vect{X}\T\vect{Y}$ da, para la pendiente,
      \begin{align*}
        \est{\beta}_1 &= -\frac{\media{X}}{\SSX}\,n\media{Y}+\frac{1}{\SSX}\sumi X_iY_i \\
        &= \frac{\sumi X_iY_i-n\media{X}\,\media{Y}}{\SSX} \\
        &= \frac{SS_{XY}}{\SSX} ,
      \end{align*}
      y, para el \emph{intercept},
      \begin{align*}
        \est{\beta}_0 &= \left(\frac1n+\frac{\media{X}^2}{\SSX}\right)n\media{Y}
          - \frac{\media{X}}{\SSX}\sumi X_iY_i \\
        &= \media{Y} + \frac{n\media{X}^2\,\media{Y}}{\SSX} - \frac{\media{X}}{\SSX}\sumi X_iY_i \\
        &= \media{Y}-\media{X}\,\frac{\sumi X_iY_i-n\media{X}\,\media{Y}}{\SSX} \\
        &= \media{Y}-\est{\beta}_1\media{X} .
      \end{align*}
    \item La fila $i$ de $\vect{X}$ es $(1,\ X_i)$ y $h_{ik}$ es el elemento $(i,k)$ de
      $\vect{H}=\vect{X}\XtXi\vect{X}\T$, luego
      \begin{align*}
        h_{ik} &= (1,\ X_i)\XtXi\begin{pmatrix} 1\\ X_k \end{pmatrix} \\
        &= \frac1n+\frac{\media{X}^2}{\SSX}-\frac{\media{X}X_k}{\SSX}
           -\frac{\media{X}X_i}{\SSX}+\frac{X_iX_k}{\SSX} \\
        &= \frac1n+\frac{\media{X}^2-\media{X}X_i-\media{X}X_k+X_iX_k}{\SSX} \\
        &= \frac1n+\frac{(X_i-\media{X})(X_k-\media{X})}{\SSX} .
      \end{align*}
  \end{enumerate}
\end{demostracion}
